\section{Propuesta de organización}
Una actividad debe definir su unidad mínima de contribución antes de abrir la convocatoria. En un reto de datos puede ser un análisis documentado; en un editatón, una mejora sustentada en fuentes; en una tesis, un resultado reproducible acompañado de interpretación. Esta unidad orienta las tareas, el apoyo necesario y la revisión final. Permite también reconocer traducción, tutoría y documentación, además de programación.

Proponemos un cierre con tres entregas: producto identificable, instrucciones de reutilización y responsable de seguimiento. Para CoAfina, una ficha por proyecto puede enlazar demostración, repositorio, versión, datos y licencia. En el editatón, el identificador de artículo y su historial permiten localizar las contribuciones. En la tesis, memoria y análisis deben enlazarse desde el depósito institucional; el depósito de 2026 de la tesis de Altuve reúne memoria, código con licencia, datos con identificador y entorno ejecutable, y muestra que esa ficha puede completarse años después de la defensa \cite{tesisaltuve}. La apertura se comprueba por componente: la dedicación CC0 de un recurso no se extiende automáticamente a datos externos o dependencias \cite{cc0}.

Retomando las definiciones y metas publicadas de Salud para todos \cite{salud}, la evaluación separa participación, producción, reutilización y continuidad. Una persona activa debe cumplir una acción definida de antemano; finalizar requiere un entregable revisado. La reutilización se documenta mediante una adaptación o ejecución identificable por otra persona. La continuidad puede expresarse mediante nuevas contribuciones o cambios de función, como pasar de participante a mentor. Registrar una revisión a tres y seis meses es una propuesta operativa para próximas actividades, con responsables y fechas establecidos desde el inicio.
